\documentclass[12pt]{article}
\usepackage[utf8]{inputenc}
\usepackage[T1]{fontenc}
\usepackage{hyperref}
\title{Kua Language Resources for K-12 Curriculum Development}
\author{Compiled by Claude Sonnet 5 (claude-sonnet-5) for Andr\'as Kornai}
\date{}
\begin{document}
\maketitle

\section{Introduction}
Kua (ISO 639-3 tyu, Glottocode kuaa1238) is an endangered Khoe-Kwadi language of eastern Botswana. Search results show that it is poorly documented. Its relation to Tsoa (ISO hio) is unclear in the sources. Joshua Project reports that it has no written form and that Bible translation is still needed. I found no school materials in or for Kua.

\section{Primary Educational Resources}
I found no primers, textbooks or curricula in Kua. The search turned up no evidence that Kua is taught in schools.

\section{Linguistic and Reference Materials}
\begin{itemize}
\item Evolex Khoe-Kwadi dictionary project, which gives comparative background on the family: \url{https://evolex.huma-num.fr/home_Khoe-Kwadi.html}
\item Wikipedia article on the related Tshwa/Kua cluster: \url{https://en.wikipedia.org/wiki/Tshwa_language}
\item Endangered Languages Project listing: \url{https://endangeredlanguages.com/node?page=904}
\end{itemize}
A search summary also mentioned a 2019 documentation project on Sasi and Kua in eastern Botswana (C. Collins, NYU). I could not confirm its portal URL, so I do not cite it.

\section{Religious and Community Resources}
Joshua Project profile (people and language status, Bible translation needed): \url{https://joshuaproject.net/index.php/languages/tyu}

\section{Digital and Audio Resources}
I found no audio or digital resources in Kua apart from the documentation project mentioned above, which I could not verify.

\section{Curriculum Development Context}
Very little exists for Kua. Any curriculum would first need community-led orthography work and documentation. The Khoe-Kwadi comparative resources above can support that work.
\end{document}